\scanpage{57}
\begin{center}
CONGRUENCE HECKE OPERATORS

by

A.O.L.ATKIN
\end{center}

1. \handwritten{\uline}{Introduction}. \quad In \handwritten{\uwave}{(2)} the au\deleted{l}thor has given an account of some results on congruences for modular forms based on real time order of discovery,ind\unsure{\deleted{e}i|\deleted{i}e}spensably aided by a computer. Here we attempt(with a certain amount of hindsight)to start at the other end,and consider what kind of relevant congruence properties might be expected \handwritten{\uline}{a} \handwritten{\uline}{priori} ; our conclusion is,that something like the results of \handwritten{\uwave}{(2)} are the most that can be hoped for.

We give no proofs in this paper. The historical aspect,and a precise account of what we can prove and what is conjecture(with overwhelming evidence)is in \handwritten{\uwave}{(2)}; detailed proofs will appear elsewhere. Our purpose here is to indicate that some \reversemarginpar\margin{\handwritten{German\\ `pee'}}\normalmarginpar``$\handwritten{\underline}{\handwritten{\mathfrak{p}}}$-adic'' theory of Hecke operators exists analogous to the exact classical theory. We sketch this classical theory of Hecke operators with a view to our analogy,so that the approach is mildly eccentric. The reader can find more satisfactory accounts in Gunning\handwritten{\uwave}{(4)},Lehner\handwritten{\uwave}{(7)}, or Rankin\handwritten{\uwave}{(8)}.

We observe finally that,while we only present our conjectures and theorems for the Fourier coefficients of modular functions(i.e.,forms of dimension $0$)and the subgroup $\handwritten{\underline}{\handwritten{\Gamma}}_{\handwritten{0}}(\handwritten{\underline}{q})$\margin{\handwritten{Cap\unsure{.|\textquotesingle}\\ gamma}} for $\handwritten{\underline}{q}$ prime,we have in fact a good deal of further evidence for forms of various dimensions(including half-integral)and other subgroups. In particular,for the reader whose interests lie in the ``romantic'' tradition of Ramanujan's famous formulae for the partition function $\handwritten{\underline}{p(n)}$:
\[
\begin{aligned}
\handwritten{\underline}{p(5n+4)}  &\handwritten{\equiv} 0 \qquad (\mathrm{mod}\ 5)  &,\\
\handwritten{\underline}{p(7n+5)}  &\handwritten{\equiv} 0 \qquad (\mathrm{mod}\ 7)  &,\\
\handwritten{\underline}{p(11n+6)} &\handwritten{\equiv} 0 \qquad (\mathrm{mod}\ 11) &,
\end{aligned}
\]
we can offer such results as
\[
\handwritten{\underline}{p(206839n+2623})\ \handwritten{\equiv}\ 0 \qquad (\mathrm{mod}\ 17).
\]
However we can only prove this by means of the multiplicative congruence properties which we consider below.

% archive sheet number: 56
